\chapter{Dữ liệu và tiền xử lý}
\section{Nguồn dữ liệu và vai trò}
Đề tài có hai nhóm dữ liệu. Nhóm thứ nhất phục vụ sentiment classification, gồm review, rating và nhãn sentiment. Nhóm thứ hai phục vụ RAG, gồm corpus review lớn hơn để truy xuất bằng FAISS. Hai nhóm dữ liệu có liên hệ qua miền ứng dụng TMĐT, nhưng vai trò khác nhau: dữ liệu sentiment dùng để huấn luyện mô hình phân loại, còn dữ liệu RAG dùng làm kho bằng chứng để trả lời câu hỏi.

\begin{table}[H]
\centering
\caption{Nguồn số liệu được dùng trong Chương 4}
\label{tab:data-source-files}
\small
\begin{tabularx}{\textwidth}{L{0.34\textwidth}Y}
\toprule
File nguồn & Vai trò \\
\midrule
\code{data_summary.json} & Số dòng dữ liệu sentiment, phân phối rating/sentiment, split và độ dài text. \\
\code{module4_index_config.json} & Cấu hình final của RAG index: 82.677 documents, embedding model, dim và index type. \\
\code{review_metadata.parquet} & Metadata final dùng để tạo lại phân phối category của RAG corpus. \\
\code{manual_precision_overall.csv} & Metric Precision@k thủ công cho retrieval baseline. \\
\bottomrule
\end{tabularx}
\vspace{0.2em}
{\footnotesize Đường dẫn đầy đủ của các file nguồn được trình bày trong phụ lục.}
\end{table}

\section{Dữ liệu sentiment}
Theo summary dữ liệu của Module 1, dữ liệu sentiment có 5.971 dòng thô. Sau làm sạch vẫn còn 5.971 dòng, sau xử lý duplicate còn 5.162 dòng. Có 46 duplicate conflict bị loại. Conflict duplicate là trường hợp nội dung hoặc bản ghi trùng nhưng nhãn/rating mâu thuẫn, nếu giữ lại sẽ gây nhiễu cho mô hình.
\begin{table}[H]\centering\caption{Tóm tắt số dòng dữ liệu sentiment}\label{tab:sentiment-counts}\begin{tabular}{lr}\toprule Giai đoạn & Số dòng \\\midrule Dữ liệu thô & 5.971 \\ Sau làm sạch & 5.971 \\ Sau xử lý duplicate & 5.162 \\ Duplicate conflict bị loại & 46 \\\bottomrule\end{tabular}\end{table}

\section{Ánh xạ rating sang sentiment}
Hệ thống ánh xạ rating sang sentiment theo quy tắc: 1--2 sao là negative, 3 sao là neutral, 4--5 sao là positive. Quy tắc này đơn giản, dễ giải thích và phù hợp với dữ liệu TMĐT có rating. Tuy nhiên, nó có hạn chế: có review 5 sao nhưng nội dung phàn nàn nhẹ vì người dùng vẫn cho shop điểm cao; cũng có review 1 sao nhưng nội dung ngắn hoặc không rõ lý do.
\begin{table}[H]\centering\caption{Phân phối sentiment 3 lớp sau dedup}\label{tab:sentiment-dist}\begin{tabular}{lrr}\toprule Sentiment & Số lượng & Tỷ lệ xấp xỉ \\\midrule Positive & 2.217 & 42,95\% \\ Negative & 2.074 & 40,18\% \\ Neutral & 871 & 16,87\% \\\bottomrule\end{tabular}\end{table}
\begin{table}[H]\centering\caption{Phân phối rating 5 lớp}\label{tab:rating-dist}\begin{tabular}{lr}\toprule Rating & Số lượng \\\midrule 1 sao & 1.079 \\ 2 sao & 995 \\ 3 sao & 871 \\ 4 sao & 794 \\ 5 sao & 1.423 \\\bottomrule\end{tabular}\end{table}

\section{Chia train/valid/test}
Dữ liệu được chia thành train, validation và test. Train dùng để học mô hình. Validation dùng để chọn hyperparameter và chọn mô hình tốt nhất. Test chỉ dùng để đánh giá cuối cùng, tránh chọn mô hình dựa trên dữ liệu test.
\begin{table}[H]\centering\caption{Phân phối train/valid/test theo sentiment}\label{tab:split-dist}\begin{tabular}{lrrr}\toprule Split & Negative & Neutral & Positive \\\midrule Train & 1.452 & 609 & 1.552 \\ Valid & 311 & 131 & 332 \\ Test & 311 & 131 & 333 \\\bottomrule\end{tabular}\end{table}

Việc giữ phân phối lớp tương đối ổn định giữa các split giúp metric đáng tin hơn. Nếu validation thiếu lớp neutral, việc chọn model theo \MacroFOne{} có thể không phản ánh đúng khả năng nhận diện neutral. Theo summary, độ dài review sentiment trung bình khoảng 13,13 token/từ. Bảng \ref{tab:text-length-stats} bổ sung thống kê tính lại từ \filepath{data/processed/sentiment_reviews.csv}.
\begin{table}[H]
\centering
\caption{Thống kê độ dài review sentiment sau xử lý}
\label{tab:text-length-stats}
\begin{tabular}{lrrrr}
\toprule
Nguồn & Min & Median & Mean & Max \\
\midrule
\filepath{data/processed/sentiment_reviews.csv} & 1 & 10,0 & 13,13 & 106 \\
\bottomrule
\end{tabular}
\end{table}
Điều này cho thấy nhiều review khá ngắn. Review ngắn có ưu điểm là tín hiệu rõ, ví dụ “giao lâu”, “hàng lỗi”. Nhưng chúng cũng thiếu ngữ cảnh, gây khó cho phân biệt neutral và mixed sentiment.

\section{Dữ liệu RAG}
Dữ liệu RAG được chuẩn hóa thành corpus phục vụ truy xuất. Theo cấu hình index final của Module 4, bản final của RAG index có 82.677 documents. Đây là con số cần ưu tiên khi trình bày bản final, vì nó phản ánh artifact đang dùng trong demo cuối.

Các summary chính của Module 4 đã được đồng bộ theo artifact final, gồm \code{review_metadata.parquet}, \code{review_faiss.index} và \code{module4_index_config.json}. Do đó, khi trình bày bản cuối, số liệu cần dùng là 82.677 documents.
\begin{table}[H]
\centering
\caption{Cấu hình RAG index bản final}
\label{tab:rag-index-config}
\small
\begin{tabularx}{\textwidth}{L{0.32\textwidth}Y}
\toprule
Thuộc tính & Giá trị \\
\midrule
Số documents & 82.677 \\
Embedding model & \code{intfloat/multilingual-e5-small} \\
Embedding dimension & 384 \\
Index type & \code{FAISS IndexFlatIP} \\
Similarity & cosine via normalized inner product \\
Passage prefix & \code{passage: } \\
Query prefix & \code{query: } \\
Index path & \code{review_faiss.index} \\
Metadata path & \code{review_metadata.parquet} \\
\bottomrule
\end{tabularx}
\end{table}


\begin{table}[H]
\centering
\caption[Phân phối ngành hàng trong RAG corpus final]{Phân phối ngành hàng trong RAG corpus final}
\label{tab:rag-category-final}
\begin{tabular}{lr}
\toprule
Ngành hàng & Số documents \\
\midrule
Fashion & 24.989 \\
Electronic & 19.645 \\
HealthBeauty & 13.739 \\
BabiesToys & 10.827 \\
HomeLifestyle & 9.798 \\
App & 1.978 \\
Cosmetic & 1.262 \\
Mobile & 439 \\
\midrule
Tổng & 82.677 \\
\bottomrule
\end{tabular}
\par\vspace{0.25em}
{\footnotesize Nguồn: \filepath{models/module4/review_metadata.parquet}.}
\end{table}

\safeimage{figures_generated/rag_category_distribution_final.png}{Phân phối ngành hàng trong RAG corpus final 82.677 documents.}{fig:rag-category-final}{0.82\textwidth}

\section{Vì sao cần xử lý duplicate và nhiễu}
Duplicate làm mô hình đánh giá sai khả năng tổng quát. Nếu một review trùng xuất hiện trong cả train và test, mô hình có thể “nhớ” review đó thay vì học quy luật. Với RAG, duplicate cũng gây vấn đề: top-k có thể trả nhiều review giống nhau, làm câu trả lời lặp và citation kém đa dạng. Review TMĐT thường chứa nhiễu như nội dung nhận xu, spam, review không liên quan, câu quá ngắn, emoji, lỗi chính tả và nội dung ngoài sản phẩm. Project giảm rủi ro bằng filtering, citation và cảnh báo evidence yếu, nhưng chưa giải quyết hoàn toàn.

\safeimage{figures/module3/sentiment_3class_distribution.png}{Phân phối sentiment 3 lớp trong dữ liệu sau xử lý.}{fig:sentiment-3class}{0.78\textwidth}
\safeimage{figures/module3/sentiment_5class_distribution.png}{Phân phối rating 5 lớp trong dữ liệu sentiment.}{fig:sentiment-5class}{0.78\textwidth}

